% ==========================================================================
% macros-cubesat.tex
% Pacotes adicionais, configurações tipográficas e macros do projeto CubeSat
% Compatibilidade: ppgmc-uesc.cls / Overleaf / TeX Live 2023-2026
% ==========================================================================

% Fonte Times New Roman para conformidade ABNT NBR 14724
\usepackage{mathptmx}

% Pacotes para tabelas de alta densidade e tipografia de precisão
\usepackage{booktabs}
\usepackage{tabularx}
\usepackage{siunitx}
\usepackage[protrusion=true,expansion=false]{microtype}

% Configuração do caminho de figuras
\graphicspath{{./04-figuras/}{./figuras/}{./}}

% --------------------------------------------------------------------------
% Configuração do siunitx (ABNT: vírgula decimal e notação científica)
% --------------------------------------------------------------------------
\sisetup{
  output-decimal-marker = {,},
  group-separator       = {.},
  group-minimum-digits  = 4,
  per-mode              = symbol,
  range-phrase          = {~a~},
  range-units           = repeat,
  exponent-product      = \times,
}

% Unidades aeroespaciais e dinâmicas customizadas
\DeclareSIUnit\Grms{\ensuremath{\mathrm{G_{rms}}}}
\DeclareSIUnit\gacc{\ensuremath{\mathrm{g}}}
\DeclareSIUnit\gravity{\ensuremath{\mathrm{G}}}
\DeclareSIUnit\Gacc{\ensuremath{\mathrm{G}}}
\DeclareSIUnit\octave{oitava}
\DeclareSIUnit\torr{Torr}

% --------------------------------------------------------------------------
% Macros de domínio e computacionais
% --------------------------------------------------------------------------
\newcommand{\code}[1]{\texttt{\detokenize{#1}}}
\newcommand{\pendente}[1]{\textbf{\textcolor{red!80!black}{[PENDENTE: #1]}}}
\newcommand{\ph}[1]{\textcolor{red!60!black}{\textit{[#1]}}}
\newcommand{\rail}{\num{8.5}\,$\times$\,\num{8.5}~\si{\milli\metre}}
\newcommand{\figplaceholder}{\fbox{\rule{0pt}{5cm}\rule{0.8\textwidth}{0pt}}}
\newcommand{\Grmsbase}{G_{\text{rms, base}}}
\newcommand{\Grmspay}{G_{\text{rms, payload}}}

% Redefinição robusta de notas para evitar estouro de margem tabular
\renewcommand{\nota}[1]{%
  \par\vspace{0.2cm}%
  \noindent{\footnotesize\textbf{Nota:} #1\par}%
}

% --------------------------------------------------------------------------
% Calibração de profundidade de numeração e Sumário (ABNT NBR 6024)
% Garante hierarquia canônica até subsubseção (nível 4) sem índices quebrados
% --------------------------------------------------------------------------
\setcounter{secnumdepth}{3}
\settocdepth{subsubsection}
\setcounter{tocdepth}{3}
\sloppy
